\documentclass[11pt,a4paper]{article}
\usepackage[utf8]{inputenc}
\usepackage[T1]{fontenc}
\usepackage[spanish,es-minimal,es-noshorthands]{babel}
\AtBeginDocument{\renewcommand{\sin}{\operatorname{sen}}\renewcommand{\arcsin}{\operatorname{arcsen}}}
\usepackage{lmodern}
\usepackage[sfdefault,lining]{FiraSans}
\usepackage[a4paper,top=2.3cm,bottom=2.2cm,left=1.8cm,right=1.8cm,headheight=1.3cm,headsep=0.4cm,footskip=1cm]{geometry}
\usepackage{amsmath,amssymb}
\usepackage{icomma}
\usepackage{graphicx}
\usepackage[table]{xcolor}
\usepackage{tikz}
\usepackage{fontawesome5}
\usepackage{ccicons}
\usepackage{fancyhdr}
\usepackage{lastpage}
\usepackage{tasks}
\usepackage{needspace}
\usepackage[most]{tcolorbox}
\usepackage{array}
\usepackage[hidelinks]{hyperref}
\AtBeginDocument{\hypersetup{pdftitle={Soluciones de la ficha de la unidad 9: Cuerpos geométricos. Áreas y volúmenes. 2º ESO},pdfsubject={Matemáticas 2º ESO},pdfkeywords={poliedros, fórmula de Euler, poliedros regulares, prisma, pirámide, cilindro, cono, esfera, área, volumen, capacidad}}}

\definecolor{navy}{HTML}{1B2748}
\definecolor{accent}{HTML}{E8603A}
\definecolor{teal}{HTML}{1F9E90}
\definecolor{soft}{HTML}{EEF1F8}
\definecolor{muted}{HTML}{5D6685}

\setlength{\parindent}{0pt}
\setlength{\parskip}{3pt}
\settasks{label=\alph*), label-format=\bfseries\color{navy}, label-width=1.3em, item-indent=1.8em, column-sep=0.6em, after-item-skip=5pt, before-skip=3pt, after-skip=4pt}

% ---------- capçalera i peu ----------
\newcommand{\llicencia}{\href{https://creativecommons.org/licenses/by-nc-sa/4.0/deed.es}{\ccbyncsa}\ \ Licencia Creative Commons Atribución-NoComercial-CompartirIgual 4.0 Internacional (CC BY-NC-SA 4.0)}
\pagestyle{fancy}
\fancyhf{}
\fancyhead[L]{\raisebox{-0.35\height}{\includegraphics[height=0.95cm]{logo-gva.png}}}
\fancyhead[C]{\small\color{muted}Soluciones de la ficha · Unidad 9: Cuerpos geométricos. Áreas y volúmenes · 2\textsuperscript{o} ESO}
\fancyhead[R]{\raisebox{-0.35\height}{\includegraphics[height=0.95cm]{logo-vermellar.png}}}
\fancyfoot[L]{\scriptsize\color{muted}\llicencia}
\fancyfoot[R]{\small\color{muted}\thepage}
\renewcommand{\headrulewidth}{0.4pt}
\renewcommand{\headrule}{\color{muted!40}\hrule height \headrulewidth}
\fancypagestyle{primera}{\fancyhf{}\renewcommand{\headrulewidth}{0pt}
  \fancyfoot[L]{\scriptsize\color{muted}\llicencia}
  \fancyfoot[R]{\small\color{muted}\thepage}}

% ---------- nivells i exercicis ----------
\newcommand{\estrelles}[1]{{\color{accent}\ifcase#1\or\faStar\faStar[regular]\faStar[regular]\or\faStar\faStar\faStar[regular]\or\faStar\faStar\faStar\fi}}
\newcounter{exer}
\newcommand{\ex}[1]{\par\needspace{4\baselineskip}\medskip\stepcounter{exer}\noindent{\color{navy}\bfseries\theexer.}\ \ {\small\estrelles{#1}}\ \ }

\newcommand{\seccio}[2]{\par\needspace{10\baselineskip}\bigskip
  \begin{tcolorbox}[colback=navy,colframe=navy,coltext=white,arc=2mm,boxrule=0pt,left=3mm,right=3mm,top=1.2mm,bottom=1.2mm]
  {\large\bfseries #1\quad #2}
  \end{tcolorbox}\vspace{-2pt}}

\newcolumntype{C}[1]{>{\centering\arraybackslash}p{#1}}
\usepackage{pifont}
% quadre de recordatori (només a la fitxa d'exercicis)
\newcommand{\recorda}[1]{\par\smallskip\begin{tcolorbox}[colback=accent!7,colframe=accent!55,arc=2mm,boxrule=0.5pt,left=2.5mm,right=2.5mm,top=1mm,bottom=1mm]\small{\color{accent}\faLightbulb[regular]}\ \textbf{Recuerda:}\ #1\end{tcolorbox}}
\newcommand{\res}[1]{\boxed{#1}}\newcommand{\rest}[1]{\fbox{#1}}

\def\fmneed{14}
% =====================================================================
% Capçalera FentMates de les fitxes: un color per a cada curs.
% Els colors són els de la pàgina inicial (LEVELS a assets/js/materials.js);
% el logo ve de _eines/logo/pdf/ (python3 _eines/logo/genera.py).
%
% Ús, al final del preàmbul de la unitat:
%   % =====================================================================
% Capçalera FentMates de les fitxes: un color per a cada curs.
% Els colors són els de la pàgina inicial (LEVELS a assets/js/materials.js);
% el logo ve de _eines/logo/pdf/ (python3 _eines/logo/genera.py).
%
% Ús, al final del preàmbul de la unitat:
%   % =====================================================================
% Capçalera FentMates de les fitxes: un color per a cada curs.
% Els colors són els de la pàgina inicial (LEVELS a assets/js/materials.js);
% el logo ve de _eines/logo/pdf/ (python3 _eines/logo/genera.py).
%
% Ús, al final del preàmbul de la unitat:
%   \input{../comu/fentmates.tex}
%   \fmcurs{1eso}
% i, al començament del document:
%   \capcalera{Matemàtiques}{1r ESO}{Fitxa d'exercicis · Unitat 2}{Divisibilitat}
% Opcional abans de \input: \def\fmneed{14} (espai mínim abans de cada secció).
% =====================================================================
\makeatletter
\def\fm@logodir{../../logo/pdf}
\providecommand{\fmneed}{10}
\def\fm@cursnom{}

% clau del curs, color i color del text damunt del color
\newcommand{\fm@curs}[3]{\@namedef{fm@color@#1}{#2}\@namedef{fm@text@#1}{#3}}
\fm@curs{1eso}{E6B84A}{navy}
\fm@curs{2eso}{3B82F6}{white}
\fm@curs{3eso}{FF7A50}{white}
\fm@curs{4eso-a}{A78BFA}{white}
\fm@curs{4eso-b}{6D28D9}{white}
\fm@curs{1batx-mi}{9B1C1C}{white}
\fm@curs{1batx-ccss}{EF6B6B}{white}
\fm@curs{2batx-mii}{1E6B47}{white}
\fm@curs{2batx-ccss}{2BB3A3}{white}

\newcommand{\fmcurs}[1]{%
  \@ifundefined{fm@color@#1}{\PackageError{fentmates}{Curs desconegut: #1}{}}{}%
  \edef\fm@tmp{\noexpand\definecolor{curs}{HTML}{\@nameuse{fm@color@#1}}%
               \noexpand\colorlet{curstext}{\@nameuse{fm@text@#1}}}\fm@tmp
  \colorlet{cursfosc}{curs!72!black}%
  \def\fm@logo{\fm@logodir/vertical-#1.pdf}%
  \def\fm@marca{\fm@logodir/marca-#1.pdf}}

% ---------- capçalera de la primera pàgina ----------
% #1 matèria (petit), #2 curs, #3 tipus de document i unitat, #4 títol
\newcommand{\capcalera}[4]{%
  \gdef\fm@cursnom{#2}%
  \par\noindent
  \begin{tcolorbox}[colback=curs,colframe=curs,boxrule=0pt,arc=2.2mm,width=3.5cm,height=2.15cm,
      valign=center,left=3mm,right=2mm,top=1mm,bottom=1mm,nobeforeafter,box align=center]
    \color{curstext}{\footnotesize #1}\par\vspace{1pt}{\fontsize{19}{21}\selectfont\bfseries #2\par}%
  \end{tcolorbox}\hspace{5mm}%
  \begin{minipage}[c]{\dimexpr\linewidth-3.5cm-5mm-2.9cm\relax}
    {\color{cursfosc}\small\bfseries\MakeUppercase{#3}}\par\vspace{2pt}
    {\color{navy}\fontsize{20}{23}\selectfont\bfseries #4\par}
  \end{minipage}\hfill
  \begin{minipage}[c]{2.6cm}\raggedleft\includegraphics[width=2.6cm]{\fm@logo}\end{minipage}\par
  \vspace{2pt}}

% ---------- barres de secció del color del curs ----------
\renewcommand{\seccio}[2]{\par\needspace{\fmneed\baselineskip}\bigskip
  \begin{tcolorbox}[colback=curs,colframe=curs,coltext=curstext,arc=2mm,boxrule=0pt,left=3mm,right=3mm,top=1.2mm,bottom=1.2mm]
  {\large\bfseries #1\quad #2}
  \end{tcolorbox}\vspace{-2pt}}

% ---------- estrelles de dificultat del color del curs (a Batxillerat no se'n posen) ----------
\@ifundefined{estrelles}{}{%
  \renewcommand{\estrelles}[1]{{\color{curs}\ifcase#1\or\faStar\faStar[regular]\faStar[regular]\or\faStar\faStar\faStar[regular]\or\faStar\faStar\faStar\fi}}}

% ---------- capçalera i peu de les pàgines ----------
% primera pàgina: logos institucionals al peu
\fancypagestyle{primera}{\fancyhf{}\renewcommand{\headrulewidth}{0pt}
  \fancyfoot[L]{\raisebox{-0.32\height}{\includegraphics[height=0.78cm]{logo-gva.png}}\hspace{5mm}%
    \raisebox{-0.32\height}{\includegraphics[height=0.78cm]{logo-vermellar.png}}}
  \fancyfoot[R]{\scriptsize\color{muted}\href{https://creativecommons.org/licenses/by-nc-sa/4.0/deed.ca}{\ccbyncsa}\ \ CC BY-NC-SA 4.0 · fentmates.org\qquad\small\thepage}}
% la resta: el símbol a l'esquerra i el curs a la dreta, amb el títol al centre com abans
\fancyhead[L]{\raisebox{-0.3\height}{\includegraphics[height=0.62cm]{\fm@marca}}}
\fancyhead[R]{\small\bfseries\color{cursfosc}\fm@cursnom}
\renewcommand{\headrule}{\color{curs!45}\hrule height \headrulewidth}
\makeatother

%   \fmcurs{1eso}
% i, al començament del document:
%   \capcalera{Matemàtiques}{1r ESO}{Fitxa d'exercicis · Unitat 2}{Divisibilitat}
% Opcional abans de \input: \def\fmneed{14} (espai mínim abans de cada secció).
% =====================================================================
\makeatletter
\def\fm@logodir{../../logo/pdf}
\providecommand{\fmneed}{10}
\def\fm@cursnom{}

% clau del curs, color i color del text damunt del color
\newcommand{\fm@curs}[3]{\@namedef{fm@color@#1}{#2}\@namedef{fm@text@#1}{#3}}
\fm@curs{1eso}{E6B84A}{navy}
\fm@curs{2eso}{3B82F6}{white}
\fm@curs{3eso}{FF7A50}{white}
\fm@curs{4eso-a}{A78BFA}{white}
\fm@curs{4eso-b}{6D28D9}{white}
\fm@curs{1batx-mi}{9B1C1C}{white}
\fm@curs{1batx-ccss}{EF6B6B}{white}
\fm@curs{2batx-mii}{1E6B47}{white}
\fm@curs{2batx-ccss}{2BB3A3}{white}

\newcommand{\fmcurs}[1]{%
  \@ifundefined{fm@color@#1}{\PackageError{fentmates}{Curs desconegut: #1}{}}{}%
  \edef\fm@tmp{\noexpand\definecolor{curs}{HTML}{\@nameuse{fm@color@#1}}%
               \noexpand\colorlet{curstext}{\@nameuse{fm@text@#1}}}\fm@tmp
  \colorlet{cursfosc}{curs!72!black}%
  \def\fm@logo{\fm@logodir/vertical-#1.pdf}%
  \def\fm@marca{\fm@logodir/marca-#1.pdf}}

% ---------- capçalera de la primera pàgina ----------
% #1 matèria (petit), #2 curs, #3 tipus de document i unitat, #4 títol
\newcommand{\capcalera}[4]{%
  \gdef\fm@cursnom{#2}%
  \par\noindent
  \begin{tcolorbox}[colback=curs,colframe=curs,boxrule=0pt,arc=2.2mm,width=3.5cm,height=2.15cm,
      valign=center,left=3mm,right=2mm,top=1mm,bottom=1mm,nobeforeafter,box align=center]
    \color{curstext}{\footnotesize #1}\par\vspace{1pt}{\fontsize{19}{21}\selectfont\bfseries #2\par}%
  \end{tcolorbox}\hspace{5mm}%
  \begin{minipage}[c]{\dimexpr\linewidth-3.5cm-5mm-2.9cm\relax}
    {\color{cursfosc}\small\bfseries\MakeUppercase{#3}}\par\vspace{2pt}
    {\color{navy}\fontsize{20}{23}\selectfont\bfseries #4\par}
  \end{minipage}\hfill
  \begin{minipage}[c]{2.6cm}\raggedleft\includegraphics[width=2.6cm]{\fm@logo}\end{minipage}\par
  \vspace{2pt}}

% ---------- barres de secció del color del curs ----------
\renewcommand{\seccio}[2]{\par\needspace{\fmneed\baselineskip}\bigskip
  \begin{tcolorbox}[colback=curs,colframe=curs,coltext=curstext,arc=2mm,boxrule=0pt,left=3mm,right=3mm,top=1.2mm,bottom=1.2mm]
  {\large\bfseries #1\quad #2}
  \end{tcolorbox}\vspace{-2pt}}

% ---------- estrelles de dificultat del color del curs (a Batxillerat no se'n posen) ----------
\@ifundefined{estrelles}{}{%
  \renewcommand{\estrelles}[1]{{\color{curs}\ifcase#1\or\faStar\faStar[regular]\faStar[regular]\or\faStar\faStar\faStar[regular]\or\faStar\faStar\faStar\fi}}}

% ---------- capçalera i peu de les pàgines ----------
% primera pàgina: logos institucionals al peu
\fancypagestyle{primera}{\fancyhf{}\renewcommand{\headrulewidth}{0pt}
  \fancyfoot[L]{\raisebox{-0.32\height}{\includegraphics[height=0.78cm]{logo-gva.png}}\hspace{5mm}%
    \raisebox{-0.32\height}{\includegraphics[height=0.78cm]{logo-vermellar.png}}}
  \fancyfoot[R]{\scriptsize\color{muted}\href{https://creativecommons.org/licenses/by-nc-sa/4.0/deed.ca}{\ccbyncsa}\ \ CC BY-NC-SA 4.0 · fentmates.org\qquad\small\thepage}}
% la resta: el símbol a l'esquerra i el curs a la dreta, amb el títol al centre com abans
\fancyhead[L]{\raisebox{-0.3\height}{\includegraphics[height=0.62cm]{\fm@marca}}}
\fancyhead[R]{\small\bfseries\color{cursfosc}\fm@cursnom}
\renewcommand{\headrule}{\color{curs!45}\hrule height \headrulewidth}
\makeatother

%   \fmcurs{1eso}
% i, al començament del document:
%   \capcalera{Matemàtiques}{1r ESO}{Fitxa d'exercicis · Unitat 2}{Divisibilitat}
% Opcional abans de \input: \def\fmneed{14} (espai mínim abans de cada secció).
% =====================================================================
\makeatletter
\def\fm@logodir{../../logo/pdf}
\providecommand{\fmneed}{10}
\def\fm@cursnom{}

% clau del curs, color i color del text damunt del color
\newcommand{\fm@curs}[3]{\@namedef{fm@color@#1}{#2}\@namedef{fm@text@#1}{#3}}
\fm@curs{1eso}{E6B84A}{navy}
\fm@curs{2eso}{3B82F6}{white}
\fm@curs{3eso}{FF7A50}{white}
\fm@curs{4eso-a}{A78BFA}{white}
\fm@curs{4eso-b}{6D28D9}{white}
\fm@curs{1batx-mi}{9B1C1C}{white}
\fm@curs{1batx-ccss}{EF6B6B}{white}
\fm@curs{2batx-mii}{1E6B47}{white}
\fm@curs{2batx-ccss}{2BB3A3}{white}

\newcommand{\fmcurs}[1]{%
  \@ifundefined{fm@color@#1}{\PackageError{fentmates}{Curs desconegut: #1}{}}{}%
  \edef\fm@tmp{\noexpand\definecolor{curs}{HTML}{\@nameuse{fm@color@#1}}%
               \noexpand\colorlet{curstext}{\@nameuse{fm@text@#1}}}\fm@tmp
  \colorlet{cursfosc}{curs!72!black}%
  \def\fm@logo{\fm@logodir/vertical-#1.pdf}%
  \def\fm@marca{\fm@logodir/marca-#1.pdf}}

% ---------- capçalera de la primera pàgina ----------
% #1 matèria (petit), #2 curs, #3 tipus de document i unitat, #4 títol
\newcommand{\capcalera}[4]{%
  \gdef\fm@cursnom{#2}%
  \par\noindent
  \begin{tcolorbox}[colback=curs,colframe=curs,boxrule=0pt,arc=2.2mm,width=3.5cm,height=2.15cm,
      valign=center,left=3mm,right=2mm,top=1mm,bottom=1mm,nobeforeafter,box align=center]
    \color{curstext}{\footnotesize #1}\par\vspace{1pt}{\fontsize{19}{21}\selectfont\bfseries #2\par}%
  \end{tcolorbox}\hspace{5mm}%
  \begin{minipage}[c]{\dimexpr\linewidth-3.5cm-5mm-2.9cm\relax}
    {\color{cursfosc}\small\bfseries\MakeUppercase{#3}}\par\vspace{2pt}
    {\color{navy}\fontsize{20}{23}\selectfont\bfseries #4\par}
  \end{minipage}\hfill
  \begin{minipage}[c]{2.6cm}\raggedleft\includegraphics[width=2.6cm]{\fm@logo}\end{minipage}\par
  \vspace{2pt}}

% ---------- barres de secció del color del curs ----------
\renewcommand{\seccio}[2]{\par\needspace{\fmneed\baselineskip}\bigskip
  \begin{tcolorbox}[colback=curs,colframe=curs,coltext=curstext,arc=2mm,boxrule=0pt,left=3mm,right=3mm,top=1.2mm,bottom=1.2mm]
  {\large\bfseries #1\quad #2}
  \end{tcolorbox}\vspace{-2pt}}

% ---------- estrelles de dificultat del color del curs (a Batxillerat no se'n posen) ----------
\@ifundefined{estrelles}{}{%
  \renewcommand{\estrelles}[1]{{\color{curs}\ifcase#1\or\faStar\faStar[regular]\faStar[regular]\or\faStar\faStar\faStar[regular]\or\faStar\faStar\faStar\fi}}}

% ---------- capçalera i peu de les pàgines ----------
% primera pàgina: logos institucionals al peu
\fancypagestyle{primera}{\fancyhf{}\renewcommand{\headrulewidth}{0pt}
  \fancyfoot[L]{\raisebox{-0.32\height}{\includegraphics[height=0.78cm]{logo-gva.png}}\hspace{5mm}%
    \raisebox{-0.32\height}{\includegraphics[height=0.78cm]{logo-vermellar.png}}}
  \fancyfoot[R]{\scriptsize\color{muted}\href{https://creativecommons.org/licenses/by-nc-sa/4.0/deed.ca}{\ccbyncsa}\ \ CC BY-NC-SA 4.0 · fentmates.org\qquad\small\thepage}}
% la resta: el símbol a l'esquerra i el curs a la dreta, amb el títol al centre com abans
\fancyhead[L]{\raisebox{-0.3\height}{\includegraphics[height=0.62cm]{\fm@marca}}}
\fancyhead[R]{\small\bfseries\color{cursfosc}\fm@cursnom}
\renewcommand{\headrule}{\color{curs!45}\hrule height \headrulewidth}
\makeatother

\fmcurs{2eso}

\begin{document}
\thispagestyle{primera}

\capcalera{Matemáticas}{2\textsuperscript{o} ESO}{Soluciones de la ficha · Unidad 9}{Cuerpos geométricos.\\ Áreas y volúmenes}
\vspace{8pt}

\begin{tcolorbox}[colback=soft,colframe=soft,arc=2mm,left=3mm,right=3mm,top=1.5mm,bottom=1.5mm]
\small Intenta resolver cada ejercicio antes de mirar su solución. Si no te sale el mismo resultado, repasa los pasos o consulta los apuntes interactivos de la unidad: encontrarás cuerpos que puedes girar y herramientas que calculan áreas y volúmenes paso a paso. Los resultados aproximados están redondeados a las centésimas.
\end{tcolorbox}

% =====================================================
\seccio{1}{Poliedros. Fórmula de Euler. Poliedros regulares}

\ex{1} \begin{center}\renewcommand{\arraystretch}{1.25}\begin{tabular}{|l|C{1.6cm}|C{1.6cm}|C{1.6cm}|l|}\hline Poliedro & $C$ & $V$ & $A$ & $C+V=A+2$ \\\hline prisma pentagonal & 7 & 10 & 15 & $17=17$ \\\hline pirámide hexagonal & 7 & 7 & 12 & $14=14$ \\\hline prisma triangular & 5 & 6 & 9 & $11=11$ \\\hline pirámide cuadrangular & 5 & 5 & 8 & $10=10$ \\\hline ortoedro & 6 & 8 & 12 & $14=14$ \\\hline\end{tabular}\end{center}

\ex{1} \begin{tasks}(3)
  \task $6+8=A+2\Rightarrow\res{A=12}$
  \task $C+12=20\Rightarrow\res{C=8}$
  \task $7+V=14\Rightarrow\res{V=7}$
\end{tasks}

\ex{1} \begin{tasks}(3)
  \task \rest{sí}
  \task no (caras de dos tipos)
  \task \rest{sí}
  \task no (caras de dos tipos)
  \task \rest{sí}
  \task no (rectángulos distintos)
\end{tasks}Un poliedro regular tiene todas las caras polígonos regulares iguales y el mismo número de caras en cada vértice.

\ex{2} \begin{center}\renewcommand{\arraystretch}{1.2}\begin{tabular}{|l|c|c|c|}\hline & Caras & Forma & Caras por vértice \\\hline Tetraedro & 4 & triángulos & 3 \\\hline Cubo & 6 & cuadrados & 3 \\\hline Octaedro & 8 & triángulos & 4 \\\hline Dodecaedro & 12 & pentágonos & 3 \\\hline Icosaedro & 20 & triángulos & 5 \\\hline\end{tabular}\end{center}

\ex{2} $3n=18\Rightarrow n=6$: base \rest{hexagonal}; \ $\res{C=8}$, \ $\res{V=12}$ \ ($8+12=18+2$).

\ex{2} $n+1=9\Rightarrow n=8$: base \rest{octogonal}; \ $\res{V=9}$, \ $\res{A=16}$ \ ($9+9=16+2$).

\ex{3} a) $8+12=20=18+2$: \rest{sí} (por ejemplo, un prisma hexagonal). \ b) $6+7=13\neq14$: \rest{no}.

\ex{3} Cada ángulo de un hexágono regular mide $120^\circ$. En un vértice deben juntarse como mínimo 3 caras, y $3\cdot120^\circ=360^\circ$: las caras quedarían \rest{planas} y no cerrarían ningún cuerpo.


% =====================================================
\seccio{2}{Prismas: elementos y área}

\ex{1} \begin{tasks}(2)
  \task $6\cdot5^2=\res{150\text{ cm}^2}$
  \task $2(7\cdot4+7\cdot3+4\cdot3)=2\cdot61=\res{122\text{ cm}^2}$
\end{tasks}

\ex{1} $D=\sqrt{36+4+9}=\sqrt{49}=\res{7\text{ cm}}$.

\ex{1} $A_L=16\cdot10=\res{160\text{ cm}^2}$; \ $A_T=160+2\cdot16=\res{192\text{ cm}^2}$.

\ex{2} Hipotenusa: $\sqrt{25+144}=13$. \ $A_L=(5+12+13)\cdot10=300$; \ $A_{\text{base}}=\frac{5\cdot12}{2}=30$. \ $A_T=300+60=\res{360\text{ cm}^2}$.

\ex{2} $A_L=48\cdot5=240$; \ $a_b=\sqrt{64-16}=\sqrt{48}\approx6{,}93$; \ $A_{\text{base}}=\frac{48\cdot\sqrt{48}}{2}\approx166{,}28$. \ $A_T\approx240+332{,}55=\res{572{,}55\text{ cm}^2}$.

\ex{2} Base $30\cdot20=600$; \ caras laterales $2\cdot30\cdot15+2\cdot20\cdot15=1500$. \ Total: $\res{2100\text{ cm}^2}$ (solo hay una base).

\ex{3} $6l^2=216\Rightarrow l^2=36\Rightarrow\res{l=6\text{ cm}}$; \ $D=\sqrt{36+36+36}=\sqrt{108}\approx\res{10{,}39\text{ cm}}$.

\ex{3} Paredes: $2(4+3{,}5)\cdot2{,}5=37{,}5$ m²; \ quitamos la puerta y la ventana: $37{,}5-1{,}6-1{,}2=34{,}7$ m². \ Coste: $34{,}7\cdot4=\res{138{,}80\text{ €}}$.


% =====================================================
\seccio{3}{Pirámides: elementos y área}

\ex{1} $a_b=8$; \ $a_p=\sqrt{36+64}=\res{10\text{ cm}}$. \ $A_T=\frac{64\cdot10}{2}+16^2=320+256=\res{576\text{ cm}^2}$.

\ex{1} $A_L=\frac{36\cdot10}{2}=\res{180\text{ cm}^2}$.

\ex{1} $A_L=\frac{24\cdot12}{2}=\res{144\text{ cm}^2}$.

\ex{2} $a_b=5$; \ $h=\sqrt{13^2-5^2}=\sqrt{144}=\res{12\text{ cm}}$. \ $A_T=\frac{40\cdot13}{2}+100=\res{360\text{ cm}^2}$.

\ex{2} $a_p=\sqrt{16+9}=\res{5\text{ cm}}$. \ La arista lateral es la hipotenusa de la apotema y medio lado: $\sqrt{25+9}=\sqrt{34}\approx\res{5{,}83\text{ cm}}$. \ $A_T=\frac{24\cdot5}{2}+36=\res{96\text{ cm}^2}$.

\ex{2} Son 4 triángulos equiláteros. Altura de cada uno: $\sqrt{16-4}=\sqrt{12}\approx3{,}46$. \ $A=4\cdot\frac{4\cdot\sqrt{12}}{2}=8\sqrt{12}\approx\res{27{,}71\text{ cm}^2}$.

\ex{3} $a_b=1$; \ $a_p=\sqrt{1{,}5^2+1^2}=\sqrt{3{,}25}\approx1{,}80$ m. \ Solo el área lateral: $A_L=\frac{8\cdot\sqrt{3{,}25}}{2}\approx\res{7{,}21\text{ m}^2}$.

\ex{3} $a_b=\sqrt{16-4}=\sqrt{12}\approx3{,}46$; \ $a_p=\sqrt{36+12}=\sqrt{48}\approx6{,}93$. \ $A_L=\frac{24\cdot\sqrt{48}}{2}\approx83{,}14$; \ $A_{\text{base}}=\frac{24\cdot\sqrt{12}}{2}\approx41{,}57$. \ $A_T\approx\res{124{,}71\text{ cm}^2}$.


% =====================================================
\seccio{4}{Cuerpos de revolución: cilindro y cono}

\ex{1} \begin{tasks}(2)
  \task \rest{cilindro}
  \task \rest{cono}
  \task \rest{esfera}
  \task \rest{dos conos} unidos por la base
\end{tasks}

\ex{1} $A_L=2\pi\cdot3\cdot8=48\pi\approx\res{150{,}80\text{ cm}^2}$; \ $A_T=48\pi+18\pi=66\pi\approx\res{207{,}35\text{ cm}^2}$.

\ex{1} $A_L=\pi\cdot6\cdot10=60\pi\approx\res{188{,}50\text{ cm}^2}$; \ $A_T=60\pi+36\pi=96\pi\approx\res{301{,}59\text{ cm}^2}$.

\ex{2} $r=8$; \ $g=\sqrt{225+64}=\sqrt{289}=\res{17\text{ cm}}$. \ $A_T=136\pi+64\pi=200\pi\approx\res{628{,}32\text{ cm}^2}$.

\ex{2} $A_L=2\pi\cdot4\cdot11=88\pi\approx\res{276{,}46\text{ cm}^2}$.

\ex{2} $h=\sqrt{25^2-7^2}=\sqrt{576}=\res{24\text{ cm}}$.

\ex{3} $2\pi r\cdot10=314{,}16\Rightarrow r=\frac{314{,}16}{20\pi}\approx\res{5\text{ cm}}$. \ $A_T=100\pi+50\pi=150\pi\approx\res{471{,}24\text{ cm}^2}$.

\ex{3} Un vaso tiene la superficie lateral y un solo círculo: $54\pi+9\pi=63\pi\approx\res{197{,}92\text{ cm}^2}$. \ 1 m² = 10\,000 cm²: $\frac{10\,000}{197{,}92}\approx50{,}5\Rightarrow\res{50\text{ vasos}}$.


% =====================================================
\seccio{5}{La esfera y su área}

\ex{1} \begin{tasks}(2)
  \task $4\pi\cdot49=196\pi\approx\res{615{,}75\text{ cm}^2}$
  \task $r=5$: \ $4\pi\cdot25=100\pi\approx\res{314{,}16\text{ cm}^2}$
\end{tasks}

\ex{1} $4\pi r^2=36\pi\Rightarrow r^2=9\Rightarrow\res{r=3\text{ cm}}$.

\ex{2} $r=3{,}25$: \ $A=4\pi\cdot3{,}25^2=42{,}25\pi\approx\res{132{,}73\text{ cm}^2}$.

\ex{2} La mitad de la esfera: $\frac{4\pi\cdot36}{2}=72\pi\approx\res{226{,}19\text{ m}^2}$.

\ex{2} $2\pi r=50\Rightarrow r=\frac{50}{2\pi}\approx\res{7{,}96\text{ cm}}$; \ $A=4\pi r^2=\frac{2500}{\pi}\approx\res{795{,}77\text{ cm}^2}$.

\ex{3} $r=5$: \ $A=100\pi\approx\res{314{,}16\text{ cm}^2}$. \ $A_{\text{cubo}}=600$: \ $\frac{314{,}16}{600}\cdot100\approx\res{52{,}36\,\%}$.

\ex{3} $A_{\text{esfera}}=4\pi\cdot16=64\pi$; \ $A_L=2\pi\cdot4\cdot8=64\pi$ \ \checkmark. \ $A_T=64\pi+32\pi=96\pi$: \ $\frac{64\pi}{96\pi}=\res{\tfrac23}$ (descubrimiento de Arquímedes).


% =====================================================
\seccio{6}{Volumen. Prismas y cilindros}

\ex{1} \begin{tasks}(3)
  \task \rest{5000 dm³}
  \task \rest{2500 cm³}
  \task \rest{7,5 dm³}
  \task \rest{350 L}
  \task \rest{3000 cm³}
  \task \rest{0,25 dm³}
\end{tasks}

\ex{1} $6^3=\res{216\text{ cm}^3}$; \ $10\cdot4\cdot5=\res{200\text{ cm}^3}$.

\ex{1} $V=\pi\cdot2^2\cdot10=40\pi\approx\res{125{,}66\text{ cm}^3}$.

\ex{2} $A_{\text{base}}=\frac{3\cdot4}{2}=6$; \ $V=6\cdot20=\res{120\text{ cm}^3}$.

\ex{2} $60\cdot30\cdot40=72\,000\text{ cm}^3=\res{72\text{ L}}$; \ $60\cdot30\cdot35=63\,000\text{ cm}^3=\res{63\text{ L}}$.

\ex{2} $V=\pi\cdot1^2\cdot3=3\pi\approx9{,}42\text{ m}^3\approx\res{9425\text{ L}}$.

\ex{3} $a_b=\sqrt{4-1}=\sqrt3\approx1{,}73$; \ $A_{\text{base}}=\frac{12\cdot\sqrt3}{2}=6\sqrt3\approx10{,}39$. \ $V=60\sqrt3\approx\res{103{,}92\text{ cm}^3}$.

\ex{3} 1 L = 1000 cm³: \ $6\cdot9{,}5\cdot h=1000\Rightarrow h=\frac{1000}{57}\approx\res{17{,}54\text{ cm}}$.


% =====================================================
\seccio{7}{Volumen de pirámides, conos y esferas}

\ex{1} $V=\frac{81\cdot10}{3}=\res{270\text{ cm}^3}$.

\ex{1} $V=\frac{\pi\cdot16\cdot9}{3}=48\pi\approx\res{150{,}80\text{ cm}^3}$.

\ex{1} $\frac43\pi\cdot27=36\pi\approx\res{113{,}10\text{ cm}^3}$; \ $\frac43\pi\cdot1000=\frac{4000\pi}{3}\approx\res{4188{,}79\text{ cm}^3}$.

\ex{2} $h=\sqrt{225-81}=\res{12\text{ cm}}$; \ $V=\frac{\pi\cdot81\cdot12}{3}=324\pi\approx\res{1017{,}88\text{ cm}^3}$.

\ex{2} $a_b=4$; \ $h=\sqrt{25-16}=\res{3\text{ cm}}$; \ $V=\frac{64\cdot3}{3}=\res{64\text{ cm}^3}$.

\ex{2} $r=12$: \ $V=\frac43\pi\cdot1728=2304\pi\approx7238{,}23\text{ cm}^3\approx\res{7{,}24\text{ L}}$.

\ex{3} $\frac43r^3=288\Rightarrow r^3=216\Rightarrow\res{r=6\text{ cm}}$; \ $A=4\pi\cdot36=144\pi\approx\res{452{,}39\text{ cm}^2}$.

\ex{3} $V_{\text{cil}}=\pi\cdot9\cdot6=54\pi$; \ $V_{\text{esfera}}=\frac43\pi\cdot27=36\pi$. \ Se derrama $\frac{36\pi}{54\pi}=\res{\tfrac23}$ del agua.


% =====================================================
\seccio{8}{Cuerpos compuestos y problemas}

\ex{1} $V=5^3+\frac{25\cdot6}{3}=125+50=\res{175\text{ cm}^3}$.

\ex{2} $V=\frac{\pi\cdot9\cdot12}{3}+\frac12\cdot\frac43\pi\cdot27=36\pi+18\pi=54\pi\approx\res{169{,}65\text{ cm}^3}$.

\ex{2} $V=\pi\cdot9\cdot8+\frac{\pi\cdot9\cdot4}{3}=72\pi+12\pi=84\pi\approx\res{263{,}89\text{ m}^3}$.

\ex{2} $V=\pi(6^2-5^2)\cdot200=2200\pi\approx\res{6911{,}50\text{ cm}^3}$.

\ex{2} $V=12^3-\pi\cdot4\cdot12=1728-48\pi\approx\res{1577{,}20\text{ cm}^3}$.

\ex{3} a) $25\cdot10+2(25+10)\cdot2=250+140=\res{390\text{ m}^2}$. \ b) $25\cdot10\cdot2=500\text{ m}^3=\res{500\,000\text{ L}}$. \ c) $\frac{500\,000}{400}=1250$ min $\approx\res{20{,}83\text{ horas}}$ (20 h 50 min).

\ex{3} Con $r=3{,}25$ y $h=6r$: \ $V_{\text{tubo}}=\pi r^2\cdot6r=6\pi r^3$; \ $V_{\text{pelotas}}=3\cdot\frac43\pi r^3=4\pi r^3$. \ Ocupan $\res{\tfrac46=\tfrac23}$ del tubo. \ Vacío: $2\pi r^3=2\pi\cdot3{,}25^3\approx\res{215{,}69\text{ cm}^3}$.

\ex{3} Volumen de las bolas: $5\cdot\frac43\pi=\frac{20\pi}{3}$. \ El agua sube $h$: \ $16\pi h=\frac{20\pi}{3}\Rightarrow h=\frac{5}{12}\approx0{,}42$ cm. \ Llega a $\res{10{,}42\text{ cm}}$.


\end{document}
